\documentclass{article}
\usepackage[english]{babel}

\usepackage[a4paper,top=2cm,bottom=2cm,left=3cm,right=3cm,marginparwidth=1.75cm]{geometry}
\usepackage{amsmath}
\usepackage{unicode-math}
\usepackage{graphicx}
\usepackage[colorlinks=true, allcolors=blue]{hyperref}
\usepackage{listings}
\usepackage{minted}
\usepackage{booktabs}
\usepackage{physics}        %braket schreibweise
\setlength{\parindent}{0pt}

\title{Physics of Complex Systems: Exercise }
\author{Leopold Schaller and Daniel Lender}

\begin{document}
\maketitle

\section*{Task 7.1: NAND and XOR}
A perceptron works the following. A input vector $\vec{x}$ is mapped to a reel number $S$ via
\begin{align}
    S = \vec{w} \vec{x} + b \text{ .}
\end{align}
Here, we have to define suitable weights $\vec{w}$ and the bias $b$. The resultning number $S$ is the evaluated with a testfunction, which is in our case the Heavyside function $\sigma$.
\subsection*{(a)}
A NAND gate evaluates an input input($x_1x_2$) and gives us 0 for the cases 00, 01, 10 and 1 for 11 case. Considering $\sigma$ as testfunction, we obtain the following conditions:

\begin{align}
x_1x_2 = 00 &: b \ge 0 \\
x_1x_2 = 10 &: w_1 + b \ge 0 \\
x_1x_2 = 01 &: w_2 + b \ge 0 \\
x_1x_2 = 11 &: w_1 + w_2 + b < 0
\end{align}
We see that $b$ has to be positive, and that $b$ has to be greater then $w_1$ and $w_2$. The last equation gives the upper limit für the weights, saying the at least on of them hast to be negative. Let us no consider symmetric inputs ($w_1 = w_2 = w$), we obtain:

\begin{align}
    b &\ge 0 \\
    b &\ge -w \\
    b &< -2w
\end{align}

For this szenario, we can choose $b$ freely and have to select the corresponding weights. If we choose $b=1$ for example, we can use $-1<w<-0.5$

For a example with different weights, the resulting values for $S$ with $b=1$,$w_1=-0.9$ and $w_2 = -0.8$ are presented:
\begin{align}
x_1x_2 = 00 &: S=1 \\
x_1x_2 = 10 &: S=0.1 \\
x_1x_2 = 01 &: S=0.2 \\
x_1x_2 = 11 &: S=-0.7
\end{align}
This outcome fulfills the NAND rule.

\subsection*{(b)}
With the same receipt, we now want to calculate the bias and weights for a XOR gate. A XOR gate responds 1 for unequal inputs 01, 10 and a 0 for equal inputs 00 and 11. This leads to the following conditions:

\begin{align}
x_1x_2 = 00 &: b < 0 \\
x_1x_2 = 10 &: w_1 + b \ge 0 \\
x_1x_2 = 01 &: w_2 + b \ge 0 \\
x_1x_2 = 11 &: w_1 + w_2 + b < 0
\end{align}

If we add up the second and third row, we obtain
\begin{align}
w_1 + w_2 + 2b  = \underset{<0}{(w_1 + w_2 + b)} + \underset{<0}{b} \ge 0
\end{align}
    But we know that the first and the second summand have to be negative. This contradiction shows, that we cannot build a XOR gate with one simple perceptron.

\end{document}

# lualatex -shell-escape Uebung6/exercise6/exercise6_Schaller_Lender.tex
